%%%%%%%%%%%%%%%%%%%%%%%%%%%%%%%%%%%%%%%%%%%%%%%%%%%%%%%%%%%%%%%%%%%%%%%%%%%%%%%%
%%  content/08-appendix.tex — اسلایدهای پشتیبان
%%
%%  این فایل *پس از* \PTclosingframe در main.tex می‌آید: اسلایدهایی که در ارائهٔ
%%  اصلی نشان داده نمی‌شوند و فقط برای پاسخ به پرسش داوران آماده‌اند.
%%
%%  همهٔ قاب‌ها [noframenumbering] دارند، پس شمارهٔ اسلاید و نوار پیشرفت دست
%%  نمی‌خورند و «۱۸ از ۱۸» در پایان ارائه همان ۱۸ می‌ماند.
%%%%%%%%%%%%%%%%%%%%%%%%%%%%%%%%%%%%%%%%%%%%%%%%%%%%%%%%%%%%%%%%%%%%%%%%%%%%%%%%
\appendix

\section{پیوست: اسلایدهای پشتیبان}

%%%%%%%%%%%%%%%%%%%%%%%%%%%%%%%%%%%%%%%%%%%%%%%%%%%%%%%%%%%%%%%%%%%%%%%%%%%%%%%%
\begin{frame}[noframenumbering]{چرا دو بازرتبه‌بند و نه یکی؟}
  \begin{center}
    \small
    \begin{tabular}{@{}lccc@{}}
      \toprule
      پیکربندی & \en{Precision@3} & هزینهٔ نسبی & تأخیر \\
      \midrule
      بدون بازرتبه‌بندی        & ۶۷٫۷٪    & ۱٫۰ & ۳٫۱ ثانیه \\
      فقط \en{Cohere}          & ۸۸٫۱٪    & ۱٫۲ & ۳٫۶ ثانیه \\
      فقط \en{Jina}            & ۸۶٫۴٪    & ۱٫۵ & ۳٫۹ ثانیه \\
      \en{Cohere} + \en{Jina}  & \hl{۹۳٫۳٪} & ۱٫۳ & ۴٫۲ ثانیه \\
      \bottomrule
    \end{tabular}
  \end{center}

  \begin{itemize}
    \item مرحلهٔ نخست \term{ارزان و گسترده} است و فهرست خام را کوچک می‌کند؛ مرحلهٔ
          دوم \term{گران و باریک} است و فقط روی ده سند اجرا می‌شود.
    \item همین ترتیب است که هزینهٔ کل را از اجرای یک‌مرحله‌ای گران‌تر پایین‌تر
          نگه می‌دارد.
  \end{itemize}
\end{frame}

%%%%%%%%%%%%%%%%%%%%%%%%%%%%%%%%%%%%%%%%%%%%%%%%%%%%%%%%%%%%%%%%%%%%%%%%%%%%%%%%
\begin{frame}[noframenumbering]{فراپارامترها و تنظیمات}
  \begin{center}
    \small
    \begin{tabular}{@{}ll@{}}
      \toprule
      پارامتر & مقدار \\
      \midrule
      تعداد پرس‌وجوی بازنویسی‌شده          & ۳ \\
      تعداد نتیجهٔ هر پرس‌وجو              & ۱۰ \\
      خروجی بازرتبه‌بند نخست               & ۱۰ سند \\
      خروجی بازرتبه‌بند دوم                & ۳ سند \\
      بیشینهٔ طول هر قطعهٔ متن             & ۱۰۰۰ نویسه \\
      دمای مدل زبانی                       & ۰٫۲ \\
      \bottomrule
    \end{tabular}
  \end{center}

  \begin{PTnote}
    دمای پایین انتخاب شده است تا مدل از سه سند داده‌شده فراتر نرود؛ همین تنظیم
    نرخ توهم‌زایی را در آزمون‌ها به صفر رساند.
  \end{PTnote}
\end{frame}

%%%%%%%%%%%%%%%%%%%%%%%%%%%%%%%%%%%%%%%%%%%%%%%%%%%%%%%%%%%%%%%%%%%%%%%%%%%%%%%%
\begin{frame}[noframenumbering]{نمونهٔ پاسخ سامانه}
  \begin{PTexample}
    \textbf{پرسش:} آخرین نسخهٔ پایدار پایتون چیست و چه زمانی منتشر شد؟

    \medskip
    \textbf{پاسخ:} نسخهٔ پایدار کنونی \en{Python 3.12} است \soft{[۱]} که در
    مهر ۱۴۰۲ منتشر شد \soft{[۲]}. مهم‌ترین تغییرش بهبود پیام‌های خطا و
    کارایی مفسر است \soft{[۱]}.
  \end{PTexample}

  \begin{itemize}
    \item هر جملهٔ پاسخ به شمارهٔ سندی گره خورده است که از آن برداشت شده؛ داور
          می‌تواند مستقیم به منبع برود.
    \item شماره‌ها به نشانی کامل صفحه در پانویس پاسخ نگاشته می‌شوند.
  \end{itemize}
\end{frame}

%%%%%%%%%%%%%%%%%%%%%%%%%%%%%%%%%%%%%%%%%%%%%%%%%%%%%%%%%%%%%%%%%%%%%%%%%%%%%%%%
\begin{frame}[noframenumbering]{ساختار دستهٔ پرسش‌های آزمون}
  \begin{center}
    \small
    \begin{tabular}{@{}lcc@{}}
      \toprule
      دستهٔ پرسش & تعداد & درستی \\
      \midrule
      رخداد تازه     & ۱۲ & ۱۰۰٪ \\
      چندبخشی        & ۱۰ & ۹۰٪  \\
      عددی           & ۸  & ۷۵٪  \\
      تعریفی         & ۱۰ & ۱۰۰٪ \\
      \midrule
      مجموع          & ۴۰ & \hl{۹۵٪} \\
      \bottomrule
    \end{tabular}
  \end{center}

  \begin{itemize}
    \item ضعیف‌ترین دسته، پرسش‌های \term{عددی} است؛ ریشهٔ خطا برداشت‌نشدن عدد از
          جدول‌های صفحهٔ وب است، نه بازیابی نادرست.
  \end{itemize}
\end{frame}
